\documentclass[11pt]{article}
% ===========================================================================
%  Shared preamble for all MPM2D worksheets — input right after \documentclass.
%  (Files beginning with "_" are skipped by mpm2d-worksheet-publish.mjs.)
% ===========================================================================
\usepackage[letterpaper,top=2.2cm,bottom=1.9cm,left=1.6cm,right=1.6cm,headheight=28pt,footskip=24pt]{geometry}
\usepackage{amsmath,amssymb}
\usepackage{cancel}
\usepackage[T1]{fontenc}
\usepackage[utf8]{inputenc}
\usepackage{lmodern}
\usepackage{xcolor}
\usepackage[most]{tcolorbox}
\usepackage{enumitem}
\usepackage{multicol}
\usepackage{fancyhdr}
\usepackage{tikz}
\usetikzlibrary{calc}
\usepackage{pgfplots}
\pgfplotsset{compat=1.18}

\definecolor{exblue}{HTML}{4A90E2}
\definecolor{exbg}{HTML}{E6F3FF}
\definecolor{qorange}{HTML}{E69138}
\definecolor{qbg}{HTML}{FFF7CC}
\definecolor{keygray}{HTML}{475569}
\definecolor{keybg}{HTML}{F1F5F9}
\definecolor{iagreen}{HTML}{14653B}
\definecolor{titleblue}{HTML}{1D4ED8}
% Rotating palette so each worked example gets its own colour.
\definecolor{ex0}{HTML}{2563A0}\definecolor{ex1}{HTML}{3B7D3B}\definecolor{ex2}{HTML}{8A5A00}
\definecolor{ex3}{HTML}{A3327A}\definecolor{ex4}{HTML}{6D28A3}\definecolor{ex5}{HTML}{B08900}
\definecolor{ex6}{HTML}{0E7490}\definecolor{ex7}{HTML}{9A3412}\definecolor{ex8}{HTML}{15803D}

\pagestyle{fancy}
\fancyhf{}
\fancyhead[L]{\footnotesize NAME: \underline{\hspace{3.4cm}}}
\fancyhead[C]{\textbf{Integration Academy}\\[-2pt]{\scriptsize Math Simplified \textperiodcentered{} Success Amplified}}
\fancyhead[R]{\footnotesize MPM2D}
\fancyfoot[L]{\scriptsize dr.merie@IntegrationAcademy.ca}
\fancyfoot[C]{\scriptsize\thepage}
\fancyfoot[R]{\scriptsize IntegrationAcademy.ca}
\renewcommand{\headrulewidth}{1.2pt}
\renewcommand{\footrulewidth}{0.4pt}
\setlength{\parindent}{0pt}
\setlength{\parskip}{4pt}

% Examples are NOT breakable, so each example + its solution stays on one page.
\newtcolorbox{example}[2]{colframe=#1,colback=#1!7,coltitle=white,colbacktitle=#1,fonttitle=\bfseries,title={#2},arc=3pt,boxrule=1pt}
\newtcolorbox{qbox}[1]{colframe=qorange,colback=qbg,coltitle=white,colbacktitle=qorange,fonttitle=\bfseries,title={#1},arc=3pt,breakable,boxrule=1pt}
\newtcolorbox{keybox}{colframe=keygray,colback=keybg,coltitle=white,colbacktitle=keygray,fonttitle=\bfseries,title={$\checkmark$\ Answer Key},arc=3pt,breakable,boxrule=1pt}
\newtcolorbox{recapbox}{colframe=keygray,colback=keybg,coltitle=white,colbacktitle=keygray,fonttitle=\bfseries,title={Question Recap},arc=3pt,breakable,boxrule=1pt}
\newtcolorbox{ideabox}{colframe=titleblue,colback=titleblue!6,coltitle=white,colbacktitle=titleblue,fonttitle=\bfseries,title={The Big Ideas},arc=3pt,breakable,boxrule=1.2pt}
\newtcolorbox{learnbox}{colframe=iagreen,colback=iagreen!5,coltitle=white,colbacktitle=iagreen,fonttitle=\bfseries,title={The Big Picture},arc=3pt,breakable,boxrule=1.2pt}
\newcommand{\lh}[1]{\par\smallskip{\bfseries\color{iagreen}#1}\par\nobreak\smallskip}

\newcommand{\soln}{\par\smallskip\textit{\textbf{Solution.}}\ }
\newcommand{\workspace}[1]{\par\vspace{3pt}\fcolorbox{black!30}{white}{\begin{minipage}[t][#1][t]{\dimexpr\linewidth-2\fboxsep\relax}\ \end{minipage}}\par}
\newcommand{\sech}[1]{\par\vspace{8pt}{\Large\bfseries\color{iagreen}#1}\par\nobreak\vspace{2pt}{\color{black!20}\hrule height1pt}\vspace{6pt}}

% A compact coordinate plot sized for an example box: \eplot{xmin}{xmax}{ymin}{ymax}{body}
\newcommand{\eplot}[5]{\begin{center}\begin{tikzpicture}
\begin{axis}[width=6.8cm,height=4.4cm,axis lines=middle,xlabel={\tiny$x$},ylabel={\tiny$y$},
grid=both,grid style={gray!16},xmin=#1,xmax=#2,ymin=#3,ymax=#4,
xtick distance=2,ytick distance=2,tick label style={font=\tiny},axis line style={-},samples=120]
#5
\end{axis}\end{tikzpicture}\end{center}}

% A sine-curve plot (degrees on x, 0..360): \splot{ymin}{ymax}{body}
\newcommand{\splot}[3]{\begin{center}\begin{tikzpicture}
\begin{axis}[width=8.6cm,height=4.6cm,axis lines=middle,xlabel={\tiny$x\,(^\circ)$},ylabel={\tiny$y$},
grid=both,grid style={gray!16},xmin=0,xmax=360,ymin=#1,ymax=#2,
xtick distance=90,ytick distance=1,tick label style={font=\tiny},axis line style={-},samples=180]
#3
\end{axis}\end{tikzpicture}\end{center}}

% A small coordinate plot: \plot{xmin}{xmax}{ymin}{ymax}{body}
\newcommand{\plot}[5]{\begin{center}\begin{tikzpicture}
\begin{axis}[width=8.4cm,height=6cm,axis lines=middle,xlabel={\scriptsize$x$},ylabel={\scriptsize$y$},
grid=both,grid style={gray!16},xmin=#1,xmax=#2,ymin=#3,ymax=#4,
xtick distance=2,ytick distance=2,tick label style={font=\tiny},axis line style={-}]
#5
\end{axis}\end{tikzpicture}\end{center}}

% Right triangle (angle theta at bottom-left, right angle at bottom-right):
%   \rtri{drawBase}{drawHeight}{oppLabel}{adjLabel}{hypLabel}
\newcommand{\rtri}[5]{\begin{center}\begin{tikzpicture}[scale=0.85]
\coordinate (A) at (0,0);\coordinate (B) at (#1,0);\coordinate (C) at (#1,#2);
\draw[thick,exblue] (A)--(B)--(C)--cycle;
\draw[black] ($(B)+(-0.32,0)$)--($(B)+(-0.32,0.32)$)--($(B)+(0,0.32)$);
\node[below] at ($(A)!0.5!(B)$){#4};
\node[right] at ($(B)!0.5!(C)$){#3};
\node[above left] at ($(A)!0.5!(C)$){#5};
\node at (0.7,0.22){$\theta$};
\end{tikzpicture}\end{center}}

% General (oblique) triangle with vertex labels and opposite-side labels:
%   \gtri{Alabel}{Blabel}{Clabel}{aSide}{bSide}{cSide}
\newcommand{\gtri}[6]{\begin{center}\begin{tikzpicture}[scale=0.85]
\coordinate (A) at (0,0);\coordinate (B) at (5,0);\coordinate (C) at (1.6,3);
\draw[thick,exblue] (A)--(B)--(C)--cycle;
\node[below left] at (A){#1};\node[below right] at (B){#2};\node[above] at (C){#3};
\node[below] at ($(A)!0.5!(B)$){#6};
\node[right] at ($(B)!0.5!(C)$){#4};
\node[left] at ($(A)!0.5!(C)$){#5};
\end{tikzpicture}\end{center}}

\fancyhead[R]{\footnotesize MCV4U \textperiodcentered{} 1.4}
\begin{document}

\begin{center}
{\LARGE\bfseries\color{titleblue}Worksheet 1.4 --- Slope of Tangent}\\[2pt]
{\small Grade 12 Calculus and Vectors (MCV4U) \textperiodcentered{} Unit 1: Rates of Change and Limits}
\end{center}

\begin{ideabox}
Now that limits are exact (Lessons 1.2--1.3), find the tangent slope at a specific point $x=a$ exactly: $m=\displaystyle\lim_{h\to0}\dfrac{f(a+h)-f(a)}{h}$. Then the tangent line is $y-f(a)=m(x-a)$.
\begin{itemize}[leftmargin=*,itemsep=2pt]
  \item $m=\displaystyle\lim_{h\to0}\dfrac{f(a+h)-f(a)}{h}$ --- expand, simplify to cancel $h$, then let $h\to0$.
  \item Tangent line at $a$: $y-f(a)=m(x-a)$.
  \item This is the exact version of Lesson 1.1's estimate.
\end{itemize}
\end{ideabox}

\begin{learnbox}
\lh{Exact, not estimated}Lesson 1.1 estimated the tangent slope with a tiny $h$. Now that limits are exact, the tangent slope at $x=a$ is $m=\displaystyle\lim_{h\to0}\dfrac{f(a+h)-f(a)}{h}$, computed precisely.
\lh{One point at a time}This limit is evaluated at a specific number $a$ — the next lesson generalizes it to every $x$ at once.
\lh{The tangent line}Once $m$ and the point $(a,f(a))$ are known, the tangent line follows from point-slope form: $y-f(a)=m(x-a)$.
\end{learnbox}

\sech{Worked Examples}

\begin{example}{ex0}{Example 1 --- Confirming the estimate}
Find the exact tangent slope of $f(x)=3x^2$ at $x=2$.\soln
\textbf{Step 1 (form the quotient):} $\dfrac{f(2+h)-f(2)}{h}=\dfrac{3(2+h)^2-12}{h}$.\\[3pt]
\textbf{Step 2 (expand):} $3(2+h)^2=3(4+4h+h^2)=12+12h+3h^2$, so the numerator is $12h+3h^2$.\\[3pt]
\textbf{Step 3 (factor and cancel $h$):} $\dfrac{12h+3h^2}{h}=12+3h$.\\[3pt]
\textbf{Step 4 (limit):} $\lim_{h\to0}(12+3h)=12$.\\[3pt]
\textbf{Conclusion:} $m=12$ exactly --- matching Lesson 1.1's estimate of $\approx12$:\eplot{-1}{3}{-2}{18}{\addplot[exblue,very thick,domain=-1:2.44]{3*x^2};\addplot[qorange,thick,domain=0.9:2.5]{12*x-12};\addplot[red,only marks,mark=*,mark size=1.6pt] coordinates {(2,12)};}
\end{example}

\begin{example}{ex1}{Example 2 --- Tangent line}
Find the tangent line to $y=x^2$ at $x=-3$.\soln
\textbf{Step 1 (find the slope):} $\dfrac{(-3+h)^2-9}{h}=\dfrac{-6h+h^2}{h}=-6+h\to-6$.\\[3pt]
\textbf{Step 2 (find the point):} $f(-3)=9$, so the point is $(-3,9)$.\\[3pt]
\textbf{Step 3 (point-slope form):} $y-9=-6(x+3)$.\\[3pt]
\textbf{Conclusion:} $y=-6x-9$:\eplot{-5}{1}{-2}{16}{\addplot[exblue,very thick,domain=-4:1]{x^2};\addplot[qorange,thick,domain=-4.2:-1]{-6*x-9};\addplot[red,only marks,mark=*,mark size=1.6pt] coordinates {(-3,9)};}
\end{example}

\begin{example}{ex2}{Example 3 --- Polynomial}
Find the exact tangent slope of $f(x)=2x^2+x$ at $x=-1$.\soln
\textbf{Step 1 (form the quotient):} $\dfrac{f(-1+h)-f(-1)}{h}=\dfrac{[2(-1+h)^2+(-1+h)]-1}{h}$.\\[3pt]
\textbf{Step 2 (expand):} $2(-1+h)^2+(-1+h)=2(1-2h+h^2)-1+h=2-4h+2h^2-1+h=1-3h+2h^2$, so the numerator is $-3h+2h^2$.\\[3pt]
\textbf{Step 3 (factor and cancel $h$):} $\dfrac{-3h+2h^2}{h}=-3+2h$.\\[3pt]
\textbf{Step 4 (limit):} $\lim_{h\to0}(-3+2h)=-3$.\\[3pt]
\textbf{Conclusion:} $m=-3$.
\end{example}

\begin{example}{ex3}{Example 4 --- Reciprocal}
Find the exact tangent slope of $f(x)=\dfrac1x$ at $x=4$.\soln
\textbf{Step 1 (form the quotient):} $\dfrac{f(4+h)-f(4)}{h}=\dfrac{\frac1{4+h}-\frac14}{h}$.\\[3pt]
\textbf{Step 2 (combine the fractions):} $\dfrac1{4+h}-\dfrac14=\dfrac{4-(4+h)}{4(4+h)}=\dfrac{-h}{4(4+h)}$.\\[3pt]
\textbf{Step 3 (divide by $h$ and cancel):} $\dfrac{-h}{4(4+h)h}=\dfrac{-1}{4(4+h)}$.\\[3pt]
\textbf{Step 4 (limit):} $\lim_{h\to0}\dfrac{-1}{4(4+h)}=\dfrac{-1}{16}$.\\[3pt]
\textbf{Conclusion:} $m=-\dfrac1{16}$.
\end{example}

\begin{example}{ex4}{Example 5 --- Radical}
Find the exact tangent slope of $f(x)=\sqrt{x}$ at $x=16$.\soln
\textbf{Step 1 (form the quotient):} $\dfrac{\sqrt{16+h}-\sqrt{16}}{h}=\dfrac{\sqrt{16+h}-4}{h}$.\\[3pt]
\textbf{Step 2 (multiply by the conjugate):} $\dfrac{(16+h)-16}{h(\sqrt{16+h}+4)}=\dfrac{h}{h(\sqrt{16+h}+4)}$.\\[3pt]
\textbf{Step 3 (cancel $h$):} $=\dfrac{1}{\sqrt{16+h}+4}$.\\[3pt]
\textbf{Step 4 (limit):} $\lim_{h\to0}\dfrac{1}{\sqrt{16+h}+4}=\dfrac{1}{4+4}$.\\[3pt]
\textbf{Conclusion:} $m=\dfrac18$.
\end{example}

\begin{example}{ex5}{Example 6 --- Polynomial}
Find the exact tangent slope of $f(x)=x^2+5x$ at $x=2$.\soln
\textbf{Step 1 (form the quotient):} $\dfrac{[(2+h)^2+5(2+h)]-14}{h}$.\\[3pt]
\textbf{Step 2 (expand):} $(2+h)^2+5(2+h)=4+4h+h^2+10+5h=14+9h+h^2$, so the numerator is $9h+h^2$.\\[3pt]
\textbf{Step 3 (factor and cancel $h$):} $\dfrac{9h+h^2}{h}=9+h$.\\[3pt]
\textbf{Step 4 (limit):} $\lim_{h\to0}(9+h)=9$.\\[3pt]
\textbf{Conclusion:} $m=9$.
\end{example}

\begin{example}{ex6}{Example 7 --- Tangent line}
Find the tangent line to $y=2x^2$ at $x=1$.\soln
\textbf{Step 1 (find the slope):} $\dfrac{2(1+h)^2-2}{h}=\dfrac{4h+2h^2}{h}=4+2h\to4$.\\[3pt]
\textbf{Step 2 (find the point):} $f(1)=2$, so the point is $(1,2)$.\\[3pt]
\textbf{Step 3 (point-slope form):} $y-2=4(x-1)$.\\[3pt]
\textbf{Conclusion:} $y=4x-2$:\eplot{-1}{3}{-2}{10}{\addplot[exblue,very thick,domain=-1:2.24]{2*x^2};\addplot[qorange,thick,domain=-1:3]{4*x-2};\addplot[red,only marks,mark=*,mark size=1.6pt] coordinates {(1,2)};}
\end{example}

\begin{example}{ex7}{Example 8 --- Linear}
Find the exact tangent slope of $f(x)=5x-3$ at $x=4$.\soln
\textbf{Step 1 (form the quotient):} $\dfrac{[5(4+h)-3]-17}{h}$.\\[3pt]
\textbf{Step 2 (expand):} $5(4+h)-3=20+5h-3=17+5h$, so the numerator is $5h$.\\[3pt]
\textbf{Step 3 (divide by $h$):} $\dfrac{5h}{h}=5$ --- no $h$ remains.\\[3pt]
\textbf{Conclusion:} $m=5$ (a line's slope is constant everywhere).
\end{example}

\begin{example}{ex8}{Example 9 --- Polynomial}
Find the exact tangent slope of $f(x)=4x^2-x$ at $x=-1$.\soln
\textbf{Step 1 (form the quotient):} $\dfrac{[4(-1+h)^2-(-1+h)]-5}{h}$.\\[3pt]
\textbf{Step 2 (expand):} $4(-1+h)^2-(-1+h)=4(1-2h+h^2)+1-h=4-8h+4h^2+1-h=5-9h+4h^2$, so the numerator is $-9h+4h^2$.\\[3pt]
\textbf{Step 3 (factor and cancel $h$):} $\dfrac{-9h+4h^2}{h}=-9+4h$.\\[3pt]
\textbf{Step 4 (limit):} $\lim_{h\to0}(-9+4h)=-9$.\\[3pt]
\textbf{Conclusion:} $m=-9$.
\end{example}

\clearpage
\sech{Your Turn --- Show Your Work}
For each question, show all steps clearly in the space provided.

\begin{qbox}{Question 1}
Find the exact tangent slope of $f(x)=x^2$ at $x=7$.\workspace{2.4cm}
\end{qbox}

\begin{qbox}{Question 2}
Find the tangent line to $y=x^2$ at $x=-4$.\workspace{2.4cm}
\end{qbox}

\begin{qbox}{Question 3}
Find the exact tangent slope of $f(x)=2x^2$ at $x=4$.\workspace{2.4cm}
\end{qbox}

\begin{qbox}{Question 4}
Find the exact tangent slope of $f(x)=\dfrac1x$ at $x=-3$.\workspace{2.4cm}
\end{qbox}

\begin{qbox}{Question 5}
Find the exact tangent slope of $f(x)=\sqrt{x}$ at $x=9$.\workspace{2.4cm}
\end{qbox}

\begin{qbox}{Question 6}
Find the exact tangent slope of $f(x)=3x^2-x$ at $x=2$.\workspace{2.4cm}
\end{qbox}

\begin{qbox}{Question 7}
Find the tangent line to $y=x^2$ at $x=-2$.\workspace{2.4cm}
\end{qbox}

\begin{qbox}{Question 8}
Find the exact tangent slope of $f(x)=5x-3$ at $x=10$.\workspace{2.4cm}
\end{qbox}

\begin{qbox}{Question 9}
Find the exact tangent slope of $f(x)=x^2+5x$ at $x=0$.\workspace{2.4cm}
\end{qbox}

\begin{qbox}{Question 10}
Find the exact tangent slope of $f(x)=4x^2-x$ at $x=1$.\workspace{2.4cm}
\end{qbox}

\begin{qbox}{Question 11}
Find the tangent line to $y=3x^2$ at $x=1$.\workspace{2.4cm}
\end{qbox}

\begin{qbox}{Question 12}
Find the exact tangent slope of $f(x)=\dfrac1x$ at $x=-1$.\workspace{2.4cm}
\end{qbox}

\begin{qbox}{Question 13 --- Challenge}
Find the exact tangent slope of $f(x)=\sqrt{x}$ at $x=1$ (rationalize).\workspace{3.5cm}
\end{qbox}

\clearpage
\sech{Question Recap \& Answer Key}

\begin{recapbox}
\begin{multicols}{2}
\small
\begin{enumerate}[leftmargin=*,itemsep=3pt]
  \item Find the exact tangent slope of $f(x)=x^2$ at $x=7$.
  \item Find the tangent line to $y=x^2$ at $x=-4$.
  \item Find the exact tangent slope of $f(x)=2x^2$ at $x=4$.
  \item Find the exact tangent slope of $f(x)=\dfrac1x$ at $x=-3$.
  \item Find the exact tangent slope of $f(x)=\sqrt{x}$ at $x=9$.
  \item Find the exact tangent slope of $f(x)=3x^2-x$ at $x=2$.
  \item Find the tangent line to $y=x^2$ at $x=-2$.
  \item Find the exact tangent slope of $f(x)=5x-3$ at $x=10$.
  \item Find the exact tangent slope of $f(x)=x^2+5x$ at $x=0$.
  \item Find the exact tangent slope of $f(x)=4x^2-x$ at $x=1$.
  \item Find the tangent line to $y=3x^2$ at $x=1$.
  \item Find the exact tangent slope of $f(x)=\dfrac1x$ at $x=-1$.
  \item Find the exact tangent slope of $f(x)=\sqrt{x}$ at $x=1$ (rationalize).
\end{enumerate}
\end{multicols}
\end{recapbox}

\begin{keybox}
\begin{multicols}{2}
\small
\begin{enumerate}[leftmargin=*,itemsep=3pt]
  \item $14$
  \item $y=-8x-16$
  \item $16$
  \item $-\tfrac19$
  \item $\tfrac16$
  \item $11$
  \item $y=-4x-4$
  \item $5$
  \item $5$
  \item $7$
  \item $y=6x-3$
  \item $-1$
  \item $\tfrac12$
\end{enumerate}
\end{multicols}
\end{keybox}

\end{document}
